\documentclass[aspectratio=169]{beamer}
\usepackage[english]{babel}
\usepackage[utf8]{inputenc}
\usepackage[T1]{fontenc}
\usepackage{lmodern}
\usepackage{listings}
\usepackage{xcolor}
\usepackage{graphicx}
\usepackage{textcomp}
\DeclareUnicodeCharacter{2014}{---}
\DeclareUnicodeCharacter{2013}{--}
\DeclareUnicodeCharacter{2212}{-}
\DeclareUnicodeCharacter{2192}{\ensuremath{\rightarrow}}
\DeclareUnicodeCharacter{00B1}{\ensuremath{\pm}}
\DeclareUnicodeCharacter{00D7}{\ensuremath{\times}}
\DeclareUnicodeCharacter{00B7}{\ensuremath{\cdot}}
\DeclareUnicodeCharacter{20AC}{EUR}
\DeclareUnicodeCharacter{2070}{\textsuperscript{0}}
\DeclareUnicodeCharacter{00B9}{\textsuperscript{1}}
\DeclareUnicodeCharacter{00B2}{\textsuperscript{2}}
\DeclareUnicodeCharacter{00B3}{\textsuperscript{3}}
\DeclareUnicodeCharacter{2074}{\textsuperscript{4}}
\DeclareUnicodeCharacter{2075}{\textsuperscript{5}}
\DeclareUnicodeCharacter{2076}{\textsuperscript{6}}
\DeclareUnicodeCharacter{2077}{\textsuperscript{7}}
\DeclareUnicodeCharacter{2078}{\textsuperscript{8}}
\DeclareUnicodeCharacter{2079}{\textsuperscript{9}}
\DeclareUnicodeCharacter{2248}{\ensuremath{\approx}}
\DeclareUnicodeCharacter{2264}{\ensuremath{\le}}

% Colors for code
\definecolor{codebg}{RGB}{245,245,245}
\definecolor{codecomment}{RGB}{0,128,0}
\definecolor{codekeyword}{RGB}{0,0,255}
\definecolor{codestring}{RGB}{163,21,21}
\definecolor{bandcolor}{RGB}{200,50,50}

\lstdefinestyle{pythonstyle}{
    language=Python,
    backgroundcolor=\color{codebg},
    basicstyle=\ttfamily\scriptsize,
    keywordstyle=\color{codekeyword},
    commentstyle=\color{codecomment},
    stringstyle=\color{codestring},
    showstringspaces=false,
    breaklines=true,
    frame=single,
    rulecolor=\color{black},
    escapeinside={(*@}{@*)},
    moredelim=**[l][\color{bandcolor}\bfseries]{\#\ ---\ NEW},
    moredelim=**[l][\color{bandcolor}\bfseries]{\#\ ---\ CHANGED},
    literate=
      {á}{{\'a}}1 {é}{{\'e}}1 {í}{{\'i}}1 {ó}{{\'o}}1 {ú}{{\'u}}1
      {Á}{{\'A}}1 {É}{{\'E}}1 {Í}{{\'I}}1 {Ó}{{\'O}}1 {Ú}{{\'U}}1
      {ñ}{{\~n}}1 {Ñ}{{\~N}}1 {ü}{{\"u}}1 {Ü}{{\"U}}1
      {¿}{{?`}}1 {¡}{{!`}}1
      {—}{{---}}1 {–}{{--}}1 {−}{{-}}1
      {±}{{\ensuremath{\pm}}}1 {×}{{\ensuremath{\times}}}1
      {→}{{\ensuremath{\rightarrow}}}1
      {°}{{\textdegree}}1
      {·}{{\ensuremath{\cdot}}}1
      {€}{{\texteuro}}1
      {≈}{{\ensuremath{\approx}}}1
      {≤}{{\ensuremath{\le}}}1
      {⁰}{{\textsuperscript{0}}}1 {¹}{{\textsuperscript{1}}}1
      {²}{{\textsuperscript{2}}}1 {³}{{\textsuperscript{3}}}1
      {⁴}{{\textsuperscript{4}}}1 {⁵}{{\textsuperscript{5}}}1
      {⁶}{{\textsuperscript{6}}}1 {⁷}{{\textsuperscript{7}}}1
      {⁸}{{\textsuperscript{8}}}1 {⁹}{{\textsuperscript{9}}}1
}

\lstset{style=pythonstyle}

% Title slide macro: \lessontitle{Lesson Number}{Title}
\newcommand{\lessontitle}[2]{
    \title{Lesson #1: #2}
    \author{}
    \institute{Tec de Monterrey}
    \begin{frame}[plain]
        \titlepage
    \end{frame}
}

% Figure frame macro: \figureframe{Title}{Image path}
\newcommand{\figureframe}[2]{
    \begin{frame}{#1}
        \begin{center}
            \includegraphics[height=0.8\textheight,width=\textwidth,keepaspectratio]{#2}
        \end{center}
    \end{frame}
}

\setbeamertemplate{navigation symbols}{}
\setbeamertemplate{footline}[frame number]
\date{}
\begin{document}
\lessontitle{04}{Crossover}
\begin{frame}{Design question}
\normalsize
Given two fixed parents, which children can an operator reach? The Planta Física case has 48 real-valued decisions: four buildings by twelve months.
\end{frame}
\begin{frame}{What to measure}
\normalsize
Reach is the support of the child distribution. Measure distinct children, new gene values, clones, and validity. Sampling illustrates; enumeration establishes support size.
\end{frame}
\begin{frame}{Baseline}
\normalsize
\begin{center}\includegraphics[width=.84\textwidth,height=.53\textheight,keepaspectratio]{../figures/offspring_01_two_parents.png}\end{center}
{\small The same fixed parents provide a comparable baseline for every operator.}
\end{frame}
\begin{frame}{One point}
\normalsize
With six distinct genes, five interior cuts and two orientations produce exactly 10 children. Each gene copies one parent: no new values appear.
\end{frame}
\begin{frame}{More cuts and uniform}
\normalsize
\begin{center}\includegraphics[width=.84\textwidth,height=.53\textheight,keepaspectratio]{../figures/offspring_03_more_cuts.png}\end{center}
{\small In this implementation: one cut 10 children, two cuts 12, three cuts 8, uniform 64.}
\end{frame}
\begin{frame}{An implementation detail}
\normalsize
Three cuts are sampled from \texttt{range(1,L-1)}, excluding the last interior boundary. The count 8 depends on this implementation; it is not a general law about more cuts.
\end{frame}
\begin{frame}{BLX-\(\alpha\)}
\normalsize
If \(m=\min(a,b)\), \(M=\max(a,b)\), and \(d=M-m\), then \(x'\sim U[m-\alpha d,M+\alpha d]\). At \(\alpha=.5\), leaving the parental interval has probability \(1/2\).
\end{frame}
\begin{frame}{Clipping is local}
\normalsize
\begin{center}\includegraphics[width=.84\textwidth,height=.53\textheight,keepaspectratio]{../figures/offspring_04_blend.png}\end{center}
{\small Blend computes new gene values and may cross bounds. Clipping only repairs gene-level bounds.}
\end{frame}
\begin{frame}{Fitness guided}
\normalsize
Evaluating proposals and keeping those that do not underperform reduces observed deterioration. Every internal trial consumes an evaluator call; an unchanged parent may be returned.
\end{frame}
\begin{frame}{Representation}
\normalsize
A real vector admits arithmetic blend. A permutation must visit each stop exactly once. Ordinary cuts can duplicate some stops and omit others.
\end{frame}
\begin{frame}{Invalid routes}
\normalsize
\begin{center}\includegraphics[width=.84\textwidth,height=.53\textheight,keepaspectratio]{../figures/offspring_06_the_divide.png}\end{center}
{\small Cut and blend do not construct valid routes in this test. Uniform crossover only recovers a valid parent.}
\end{frame}
\begin{frame}{Ordered crossover}
\normalsize
OX1 copies a segment and fills gaps in the other parent's relative order. It preserves membership and uniqueness for permutations; it may be inert on real-valued vectors.
\end{frame}
\begin{frame}{Two families}
\normalsize
\begin{center}\includegraphics[width=.84\textwidth,height=.53\textheight,keepaspectratio]{../figures/offspring_07_ordered.png}\end{center}
{\small OX1 works for routes; arithmetic crossover works for real values. Choosing the family is part of the model.}
\end{frame}
\begin{frame}{Transfer to the challenge}
\normalsize
For the 48 proportions, a real-valued crossover proposes a child. The shared evaluator checks gene bounds, annual quota, potable minimum, and monthly variation. Gene clipping does not prove global feasibility.
\end{frame}
\begin{frame}{Design record}
\normalsize
In teams of three: choose family, parameter, repair, and diversity measure. The reviewer explains the method; the verifier counts evaluations and checks feasibility.
\end{frame}
\begin{frame}{Interpretation}
\normalsize
Validity rate is not fitness. Compare operators under the same budget, eight seeds, and common repair policy. These results are synthetic, not Planta Física confirmed parameters.
\end{frame}
\end{document}
